\section{Materials and Methods}
\label{ch:methods}

\subsection{Introduction}
\label{sec:m_intro}

This chapter describes how the tomato ripening detection and sorting system was built. The work has two parts: a computer vision part (collecting images, training and testing a detection model) and an engineering part (building the conveyor, its electronics and its software). Together they form a conveyor prototype that detects a tomato on a moving belt, classifies its ripening stage or defect condition, estimates its remaining shelf life, and diverts it into the correct output stream.

The chapter follows the order the work was done in: materials, dataset preparation, model development, real-time detection, hardware and software development, integration, shelf-life observation, and the experiments and tests used to evaluate the system. Results are reported in Chapter~\ref{ch:results}; raw records and full component details are in the appendices.

\subsection{Overall Research Methodology}
\label{sec:m_overall}

The project was developed iteratively rather than in one straight pass, so the dataset, annotation rules and model could be revised whenever evaluation exposed a problem. The stages, in order, were: problem definition; sample collection and image acquisition; dataset preparation, annotation and splitting; preprocessing and augmentation; model training and testing; correcting the dataset and retraining where needed; hardware design and fabrication; software development; hardware and software integration; real-time detection and sorting; and system testing. This workflow, including the correction loop, is shown in Figure~\ref{fig:methodology}.

\begin{figure}[H]
\centering
\includegraphics[width=0.5\linewidth]{figures/overallmetho.png}
\caption{Overall research methodology, from problem definition to system level evaluation}
\label{fig:methodology}
\end{figure}

\subsection{Research Materials}
\label{sec:m_materials}

\subsubsection{Tomato Samples}
\label{sec:m_samples}

Padma variety tomatoes were obtained from a commercial collecting centre in Bandarawela, Badulla District, which receives fruit from several farms in the region. Buying from a multi-farm centre rather than one farm gives natural variation in size, texture, maturity and handling damage. Fruits were sorted by eye at capture into the five categories defined in Section~\ref{sec:m_classes}. Materials used are listed in Table~\ref{tab:materials}.

\begin{table}[H]
\caption{Materials and samples used in the research}
\label{tab:materials}
\centering
\begin{tabularx}{\textwidth}{p{3.1cm} X p{2.7cm}}
\toprule
\textbf{Material} & \textbf{Description} & \textbf{Purpose} \\
\midrule
Padma tomato fruit &
Fresh fruit at four ripening stages, and fruit with visible cracking, bruising or rot, from a commercial collecting centre, Bandarawela &
Image dataset, shelf-life observation and conveyor testing \\
\addlinespace
Publicly sourced defect images &
Third-party defect images used in an earlier dataset version, taken under uncontrolled background and lighting &
Removed during quality control (Section~\ref{sec:annotation_qc}) \\
\bottomrule
\end{tabularx}
\end{table}

A total of 8{,}000 images were photographed over a 3-day collection period.

\TODO{inclusion/exclusion criteria applied per class -- state what made a captured image acceptable or caused it to be rejected}

\subsubsection{Hardware Components}
\label{sec:m_hardware}

The prototype is a belt conveyor with a fixed overhead camera and four servo-actuated sorting gates. Its main components are listed in Table~\ref{tab:hardware}; the full pin assignment and wiring diagrams are in Appendix~C.

\begin{table}[H]
\caption{Principal hardware components of the conveyor sorting prototype}
\label{tab:hardware}
\centering
\begin{tabularx}{\textwidth}{p{3.0cm} X c}
\toprule
\textbf{Component} & \textbf{Model / specification} & \textbf{Qty} \\
\midrule
Microcontroller & ESP32 development board (DOIT DEVKIT V1); runs the sorting firmware & 1 \\
Drive motor & NEMA 23 stepper motor & 1 \\
Motor driver & DM542 step-and-direction driver & 1 \\
Power supplies & 24\,V 5\,A driver rail; 12\,V 2\,A logic rail & 1 each \\
Gate actuators & MG995 servo with printed diverter blade, one per ripening class & 4 \\
Presence sensors & Infrared reflective sensor, one at each gate & 4 \\
Power transmission & Motor pulley, roller pulley and timing belt (3:1 reduction) & 1 set \\
Frame and mounts & 25\,mm box-section frame, bearing assemblies, printed brackets, adjustable camera mount & 1 set \\
Speed control & Potentiometer read by the microcontroller & 1 \\
Camera & \acrshort{USB} webcam, 960\,$\times$\,720 pixels, 30~\acrshort{FPS} & 1 \\
Host computer & Workstation with NVIDIA GeForce RTX 3070 Ti, 8{,}192~\acrshort{MiB} \acrshort{VRAM} & 1 \\
\bottomrule
\end{tabularx}
\end{table}

\subsubsection{Software and Technologies}
\label{sec:m_software}

The software environment is listed in Table~\ref{tab:software}.

\begin{table}[H]
\caption{Software, frameworks and libraries used in the research}
\label{tab:software}
\centering
\begin{tabularx}{\textwidth}{p{4.3cm} X}
\toprule
\textbf{Software} & \textbf{Purpose} \\
\midrule
Python, PyTorch/torchvision, Ultralytics (YOLOv8) & Model training, validation and inference \\
OpenCV & Camera capture, frame processing, box rendering \\
Albumentations & Offline bounding-box-aware augmentation \\
Streamlit and Plotly & Operator dashboard and charts \\
SQLite & Local detection and indicator database \\
Roboflow & Annotation and versioned dataset export \\
PlatformIO, Arduino framework, ESP32Servo & Firmware build, upload and servo control \\
SolidWorks and Onshape & Three-dimensional design of the prototype \\
\bottomrule
\end{tabularx}
\end{table}

\subsection{Sampling Method and Sample Collection}
\label{sec:m_sampling}

Images were collected purposively rather than randomly: fruit was deliberately photographed across all five target categories so that even the comparatively uncommon defect condition was well represented \citep{etikan2016}.

Collection began in January 2026. All dataset images were photographed with the fruit against a white background inside a lightbox under \acrshort{LED} illumination, keeping background and lighting the same across sessions and classes. This mattered: as Section~\ref{sec:annotation_qc} describes, a class photographed under different conditions from the rest let the model learn the lighting instead of the fruit.

\subsection{Image Acquisition}
\label{sec:m_acquisition}

Live deployment uses a separate camera path from dataset capture. For real-time inference the system uses the \acrshort{USB} webcam mounted above the conveyor, at 960\,$\times$\,720 pixels and 30~\acrshort{FPS}. After start-up, the camera's automatic exposure and white balance are given about 1.5~seconds to settle before any frame is passed to the model, since frames grabbed during that window can have colour casts strong enough to change the predicted class.

\TODO{camera model and file format used for dataset capture; camera-to-fruit distance and viewing angle; whether multiple orientations of the same fruit were captured}

\begin{figure}[H]
\centering
\TODO{Figure 3.2 -- lightbox capture setup and the conveyor-mounted camera arrangement}
\caption{Image acquisition setup for dataset capture and the conveyor mounted camera arrangement}
\label{fig:acquisition}
\end{figure}

\subsection{Dataset Description}
\label{sec:m_dataset}

The dataset was built and versioned in Roboflow \citep{dwyer2024} and exported in the YOLOv8 label format. Its composition is summarised in Table~\ref{tab:dataset} and the per-class counts in Table~\ref{tab:instances}. Instance counts are higher than image counts because one image can contain more than one fruit. The images added to boost the defect class (Section~\ref{sec:m_preprocessing}) are always listed separately, so synthetic and captured data are never mixed together in a single figure.

\begin{table}[H]
\caption{Composition of the dataset used for model development}
\label{tab:dataset}
\centering
\begin{tabularx}{\textwidth}{p{6.2cm} X}
\toprule
\textbf{Property} & \textbf{Value} \\
\midrule
Number of classes & 5 (breaker, defect, green, red, turning) \\
Training images & 5{,}498 captured, plus 668 offline-augmented defect images \\
Validation images & 1{,}574 \\
Test images (held out) & 777 \\
Total distinct captured images & 7{,}849 \\
Annotation format & YOLO normalised bounding box: class index, $x$ centre, $y$ centre, width, height \\
Model input resolution & 640\,$\times$\,640 pixels after resizing \\
\bottomrule
\end{tabularx}
\end{table}

\begin{table}[H]
\caption{Per class annotated object instance counts by dataset split}
\label{tab:instances}
\centering
\begin{tabular}{lrrrrr}
\toprule
\textbf{Class} & \textbf{Train} & \textbf{Train} & \textbf{Validation} & \textbf{Test} & \textbf{Total} \\
 & \textbf{(captured)} & \textbf{(augmented)} & & & \\
\midrule
breaker & 1{,}137 & 0   & 326 & 191 & 1{,}654 \\
defect  & 448   & 896 & 137 & 66  & 1{,}547 \\
green   & 1{,}380 & 0   & 429 & 191 & 2{,}000 \\
red     & 1{,}434 & 0   & 395 & 190 & 2{,}019 \\
turning & 1{,}138 & 0   & 309 & 149 & 1{,}596 \\
\midrule
\textbf{Total} & \textbf{5{,}537} & \textbf{896} & \textbf{1{,}596} & \textbf{787} & \textbf{8{,}816} \\
\bottomrule
\end{tabular}
\end{table}

The defect class has only 66 test instances, against 149--191 for each ripening class, so it is treated with more caution in the evaluation approach of Section~\ref{sec:m_evaluation}.

\begin{figure}[H]
\centering
\TODO{Figure 3.3 -- representative captured images for each of the five target classes}
\caption{Representative sample images for each of the five target classes}
\label{fig:classes}
\end{figure}

\subsection{Ripening Stage Classification Criteria}
\label{sec:m_classes}

Five classes were defined: four ripening stages and one quality class. Defect is a separate class rather than a property of a ripening class, because damaged fruit is handled the same way no matter what colour it is.

The four ripening stages come from merging adjacent pairs of the six standard \acrshort{USDA} maturity stages \citep{usda1991}, the same approach \citet{abekoon2024} used for this variety. Four classes map directly onto the four sorting gates, and the two narrowest standard bands sit exactly where a camera struggles most to judge red surface percentage, so merging them gives boundaries the model can separate reliably on a moving belt. The correspondence is in Table~\ref{tab:merge}, and the resulting criteria in Table~\ref{tab:ripeness}.

\begin{table}[H]
\caption{Correspondence between the six standard maturity stages and the four operational classes used in this research}
\label{tab:merge}
\centering
\begin{tabular}{llll}
\toprule
\textbf{Standard stage} & \textbf{Standard red surface} & \textbf{Operational class} & \textbf{Operational red surface} \\
\midrule
Green      & 0\,\%      & Green   & 0\,\% \\
\midrule
Breaker    & $\leq$ 10\,\%   & \multirow{2}{*}{Breaker} & \multirow{2}{*}{$<$ 30\,\%} \\
Turning    & 10--30\,\% &         & \\
\midrule
Pink       & 30--60\,\% & \multirow{2}{*}{Turning} & \multirow{2}{*}{30--80\,\%} \\
Light red  & 60--90\,\% &         & \\
\midrule
Red        & $>$ 90\,\%  & Red     & $>$ 80\,\% \\
\bottomrule
\end{tabular}

\vspace{2mm}
{\footnotesize Standard stages and boundaries from \citet{usda1991}. The operational upper boundary of the turning class is set at 80\,\% rather than 90\,\% so that fully coloured fruit falls into the red class under belt lighting.}
\end{table}

\begin{table}[H]
\caption{Ripening stage and quality classes used in this research}
\label{tab:ripeness}
\centering
\begin{tabularx}{\textwidth}{p{1.9cm} p{3.3cm} X}
\toprule
\textbf{Class} & \textbf{Stage} & \textbf{Identification characteristics} \\
\midrule
Green   & Unripe (Stage 1)             & Entirely green skin, no red or orange colouration \\
Breaker & Onset of ripening (Stage 2)  & Green-to-yellow transition, less than 30\,\% red or pink \\
Turning & Mid ripening (Stage 3)       & 30--80\,\% of the surface red, mixed green/orange/red patches \\
Red     & Fully ripe (Stage 4)         & More than 80\,\% of the surface red, uniform ripe colouration \\
Defect  & Rejected quality             & Visible cracking, bruising or rot; can occur at any stage, receives no ripening label \\
\bottomrule
\end{tabularx}
\end{table}

\subsection{Image Annotation}
\label{sec:m_annotation}

All images were annotated in Roboflow \citep{dwyer2024} and exported in the YOLOv8 format, where each image has a label file with one line per annotated fruit: a class index followed by the box's position and size.

Every fruit visible in an image was given its own box and a single class label. For the defect class the box still covers the whole fruit, not just the blemish, and is labelled defect regardless of the fruit's underlying colour. Keeping the box style the same across classes matters, because the model also learns the typical size of the box for each class along with its appearance.

\begin{figure}[H]
\centering
\TODO{Figure 3.4 -- annotated image showing whole-fruit bounding boxes on mixed classes}
\caption{Example of whole fruit bounding box annotation for a multiple fruit image with mixed classes}
\label{fig:annotation}
\end{figure}

\subsubsection{Annotation Quality Control and Correction}
\label{sec:annotation_qc}

Annotation went through a review and correction cycle. Two corrections shaped the final dataset.

\textbf{Inconsistent box style.} Defect images taken from a public dataset had been boxed around only the blemished area, while the project's own images boxed the whole fruit. Every such image was re-annotated to the whole-fruit style so all five classes were consistent.

\textbf{Lighting acting as a false class cue.} Even after the box style was fixed, live camera tests showed healthy breaker, turning and red fruit being wrongly classified as defect. The cause was that most defect images still came from the public dataset, with different backgrounds and lighting from the other four classes, which had all been shot inside the lightbox. The model had partly learned the lighting, not the damage, as the sign of a defect. The fix was to drop the mismatched public images, capture and label new defect images under the same lightbox conditions as the other classes, and use the offline augmentation described in Section~\ref{sec:m_preprocessing} to make up for the smaller number of real defect photos that remained. The dataset in Section~\ref{sec:m_dataset} is this corrected version.

\TODO{report the annotation reliability check: sample size re-annotated independently and the resulting agreement figure, since class boundaries are judged by eye rather than measured}

The detailed annotation procedure is given in Appendix~A.

\subsection{Image Preprocessing and Augmentation}
\label{sec:m_preprocessing}

All images are resized to 640\,$\times$\,640 pixels for both training and inference, the standard input size for the architecture.

During training, images are augmented on the fly: colour jitter, rotation, translation, scaling, horizontal flipping (never vertical, since fruit on a belt is never upside down), mosaic, mixup and random erasing. This kind of augmentation is standard practice for improving generalisation from limited data \citep{shorten2019}, used here because live conveyor conditions differ from the lightbox conditions the dataset was captured in.

After the corrections in Section~\ref{sec:annotation_qc}, the defect class had far fewer real training images than the ripening classes. Its training images were therefore supplemented offline using Albumentations \citep{buslaev2020}: flip, rotation, scale, brightness and hue changes produced two augmented copies of each defect training image. These extra images were used only for training; the validation and test sets contain captured images only.

An automated check runs before every training run and stops it if something is structurally wrong: it confirms the class list, checks a test split exists, counts images and instances per class, flags missing or empty label files and duplicate images, and flags any class whose boxes are unusually small compared to the others. That last check was added after the box-style problem above, so an inconsistent box convention cannot slip into training again unnoticed.

\TODO{confirm whether images of one physical fruit were confined to a single split; several photos were taken per fruit, and if the same fruit ended up split across train/validation/test, that is a limitation on the reported metrics}

\begin{figure}[H]
\centering
\TODO{Figure 3.5 -- preprocessing and augmentation pipeline, showing the offline defect branch separately}
\caption{Image preprocessing and augmentation pipeline, including the offline defect class branch}
\label{fig:preprocessing}
\end{figure}

\subsection{Dataset Splitting}
\label{sec:m_splitting}

The dataset was split into training, validation and test sets at export time, in the proportions in Table~\ref{tab:split}, a standard 70:20:10 split. Training fits the model, validation is checked after every epoch to choose the best checkpoint, and the test set is used only once, at the end, to measure final performance.

\begin{table}[H]
\caption{Dataset split used for model development}
\label{tab:split}
\centering
\begin{tabularx}{\textwidth}{l r r X}
\toprule
\textbf{Subset} & \textbf{Images} & \textbf{Percentage} & \textbf{Role} \\
\midrule
Training   & 5{,}498 & 70.0\,\% & Fitting model parameters; supplemented with 668 augmented defect images \\
Validation & 1{,}574 & 20.1\,\% & Monitoring generalisation, early stopping, checkpoint selection \\
Test       & 777   & 9.9\,\%  & Held out from training; used once for final evaluation \\
\midrule
\textbf{Total} & \textbf{7{,}849} & \textbf{100\,\%} & \\
\bottomrule
\end{tabularx}
\end{table}

\subsection{Model Development}
\label{sec:m_model}

The system uses YOLOv8, an object detection architecture that finds and classifies objects in a single pass over the image, without a separate step to first propose regions of interest \citep{redmon2016, terven2023}. The medium-sized YOLOv8m variant was used, starting from publicly released pre-trained weights and then fine-tuned on the project's own dataset.

This architecture fits the task well: a single camera frame above the conveyor can contain more than one fruit at different ripening stages, and the sorting logic needs a box, a class and a confidence score for each one, fast enough to keep up with the belt. Published work supports this choice for tomato detection \citep{safira2026}, and the medium variant balances accuracy against the memory available on the host \acrshort{GPU}. No separate classification step is used afterwards -- the sorting decision is the detector's own prediction, combined across frames as described in Section~\ref{sec:m_realtime}.

A second detector (YOLO26m) was also trained on the same data, so the architecture choice could be checked against this project's own evidence rather than published results alone; this comparison is Experiment~2 in Section~\ref{sec:m_experiments}.

\begin{figure}[H]
\centering
\TODO{Figure 3.6 -- single stage detection workflow of the YOLOv8 architecture: input frame, backbone and neck feature extraction, detection head, boxes with class and confidence}
\caption{Simplified single stage detection workflow of the YOLOv8 architecture}
\label{fig:yolo}
\end{figure}

\subsection{Model Training}
\label{sec:m_training}

Training was done locally on a workstation with an NVIDIA GeForce RTX 3070 Ti \acrshort{GPU} (8{,}192~\acrshort{MiB} \acrshort{VRAM}). Because this GPU has limited memory, the training script automatically picks a safe batch size and checks free memory before each run starts. The configuration used is in Table~\ref{tab:training}; the full configuration file is in Appendix~B.

\begin{table}[H]
\caption{Training configuration used for model development}
\label{tab:training}
\centering
\begin{tabularx}{\textwidth}{p{5.4cm} X}
\toprule
\textbf{Parameter} & \textbf{Value} \\
\midrule
Base architecture & YOLOv8m, initialised from pre-trained weights \\
Training data & 5{,}498 captured images and 668 offline-augmented defect images \\
Maximum epochs and early stopping & 150 epochs; halts after 25 epochs without validation improvement \\
Batch size and input size & Automatic selection; 640\,$\times$\,640 pixels \\
Initial learning rate & 0.01 \\
Augmentation & Colour jitter; rotation, translation, scale; horizontal flip only; mosaic, mixup and random erasing \\
\bottomrule
\end{tabularx}
\end{table}

Training runs in five steps: check the framework, device and free memory are ready; run the dataset audit from Section~\ref{sec:m_preprocessing}, which stops the run if something is wrong; train using the settings above, checking performance on the validation set after every epoch and saving the best checkpoint; stop early once performance stops improving, and keep the best checkpoint as the final model; then run one separate evaluation pass of that checkpoint on the held-out test set.

\subsection{Validation and Evaluation Procedure}
\label{sec:m_evaluation}

The validation set was checked after every epoch during training to catch overfitting early, using a fixed stopping rule rather than manual judgement. The test set was evaluated only once, so the figures reported in Chapter~\ref{ch:results} are an honest estimate rather than numbers the model was tuned towards.

Performance is measured with the standard object detection metrics \citep{padilla2020}: precision (how many of the model's positive predictions were correct), recall (how many of the real fruits of a class were actually found), an $F_1$ score that balances the two (used to pick the confidence threshold used at deployment), and mean average precision at a looser overlap threshold (mAP@50) and a stricter one that also rewards a tightly placed box (mAP@50--95). A confusion matrix is also produced to show which classes get mixed up with each other; as Section~\ref{sec:annotation_qc} shows, this diagnostic caught a problem that the summary numbers alone did not.

Because the test set is limited in size -- and the defect class especially, with only 66 instances -- a confidence interval is calculated for each reported metric by resampling the test set, so a difference between two models is only treated as real when the intervals do not overlap.

System-level figures -- fruits processed, class distribution, defect ratio, throughput and average shelf life -- are simply totalled directly from the detection log.

\subsection{Real Time Detection Process}
\label{sec:m_realtime}

At run time the system works through the following steps for every frame.

\begin{enumerate}
    \item \textbf{Frame capture.} A live frame is captured from the webcam at 960\,$\times$\,720 pixels and 30~\acrshort{FPS}.
    \item \textbf{Preprocessing.} The frame is resized to the 640\,$\times$\,640 pixel model input size.
    \item \textbf{Inference.} The frame is passed through the trained model at a confidence threshold of 0.45, the point where precision and recall are best balanced.
    \item \textbf{Non-tomato object filtering.} A second, general-purpose detector trained on everyday objects \citep{lin2014} checks the same frame for things like bottles, cups, phones and hands. Any tomato detection that overlaps one of these is discarded, since the tomato model itself has no way to recognise an object that isn't a tomato.
    \item \textbf{Tracking and vote accumulation.} A simple tracker follows one fruit at a time across frames and records a class vote from each frame it appears in. A track ends after 0.6~seconds with no detection, and needs at least two frames to count.
    \item \textbf{Classification finalisation.} When a track ends, the final class is the one with the most weight once each frame's vote is scaled by how confident the model was in it. Deciding from several frames rather than one makes the result more robust to a fruit briefly rotating or catching a bright reflection.
\end{enumerate}

The defect class has one extra, deliberately strict rule: it is only chosen as the final answer if it was voted in a clear majority of the fruit's frames and with consistently high confidence. This is a deliberate trade-off -- wrongly rejecting a good fruit loses money permanently, while letting a borderline one through can still be caught later -- and it also gives a second safeguard against the lighting confound described in Section~\ref{sec:annotation_qc}.

The final class, its confidence and its shelf-life value (Section~\ref{sec:m_shelflife}) are then written to the database and, if the hardware is connected, turned into a sorting command as described in Section~\ref{sec:m_integration}.

\begin{figure}[H]
\centering
\TODO{Figure 3.7 -- flowchart of the real time detection and classification process}
\caption{Real time tomato detection and classification process}
\label{fig:detection}
\end{figure}

\subsection{Hardware Prototype Development}
\label{sec:m_prototype}

The conveyor frame was built from 25\,mm metal box-section bar, to dimensions first checked in a \acrfull{CAD} model. The belt is driven by a NEMA~23 stepper motor through a 3:1 pulley-and-belt reduction, supported on bearing assemblies. The motor is commanded through a DM542 driver, with belt speed set by a potentiometer read by the microcontroller. An adjustable camera mount lets position and angle be tuned on the finished rig.

Four infrared sensors and four servo gates are fitted, one pair per ripening class. The defect class deliberately has no gate: a fruit classified as defective gets no command and simply carries on to the end of the belt into a single reject stream.

The order of the gates matters for the control logic in Section~\ref{sec:m_integration}: going downstream from the camera, the gates are encountered as red, turning, breaker, green. Each gate has its own sensor just before its diverter blade, and its own dedicated servo and sensor wiring, with the class-to-gate mapping fixed in firmware. Each gate rests at 180\textdegree{} and swings to 110\textdegree{} to divert a fruit, then returns automatically.

A single ESP32 board controls everything, powered from a 24\,V supply on the motor side and a separate 12\,V supply on the logic side, so motor and servo noise cannot disturb the control electronics. The pin assignment and wiring diagram are in Appendix~C.

\begin{figure}[H]
\centering
\TODO{Figure 3.8 -- CAD model and photographs of the assembled conveyor prototype}
\caption{Computer aided design model and assembled hardware prototype}
\label{fig:prototype}
\end{figure}

\subsection{Software System Development}
\label{sec:m_softwaredev}

The software is a Python application built around a web-based operator dashboard, backed by a detection module and a local database. The dashboard shows a live camera feed with boxes drawn on it; the current tracking state and the most recent fruit's class and confidence; an indicator panel with fruits processed, throughput, defect ratio and average shelf life over a chosen time window; a class-distribution chart; and a table of recent detections.

The model is loaded once per session, with a menu to pick between trained checkpoints, and every frame runs at a confidence threshold the operator can adjust (0.45 by default). Every finalised classification is saved to the database with its confidence, box, timestamp and shelf-life value, and the indicator panel reads from that same database, so what is shown on screen always matches what is stored.

\begin{figure}[H]
\centering
\TODO{Figure 3.9 -- software architecture diagram and dashboard screenshot}
\caption{Software architecture and operator dashboard of the detection and sorting application}
\label{fig:software}
\end{figure}

\subsection{Hardware and Software Integration}
\label{sec:m_integration}

The dashboard talks to the controller over Bluetooth. On classification, the host sends a short command: one of four gate commands, or a no-gate command for a defective fruit. Each command has a sequence number and is resent until acknowledged, up to a fixed number of tries. If the controller sees a command it has already acted on, it acknowledges it again but does not act on it twice. Both parts matter -- a lost command drops a fruit from the queue, and a lost acknowledgement duplicates it -- and either fault throws off every gate assignment that comes after it.

On the controller side, each gate keeps its own queue of pending classifications, in the order the fruit will reach them along the belt. When a gate's sensor detects an arriving fruit, the item at the front of that gate's queue is checked: if it matches the class the gate handles, the gate opens; if not, the item moves on to the next gate's queue to be checked again there. This keeps every classification lined up with the physical fruit it belongs to, even with several fruits on the belt at once.

Each gate is released by its own sensor rather than by a time calculated from belt speed and distance (the simpler approach outlined in Section~\ref{sec:actuation}), which removes the main source of error that would otherwise grow worse for gates further down the belt. A small delay, scaled to belt speed, compensates for the short gap between the sensor and the diverter blade itself; the gate then stays open briefly before returning to rest.

\begin{figure}[H]
\centering
\TODO{Figure 3.10 -- communication flow from classification through the staged queue to gate servo actuation}
\caption{Hardware and software integration and communication flow}
\label{fig:integration}
\end{figure}

\subsection{Shelf Life Observation Procedure}
\label{sec:m_shelflife}

Shelf life was measured directly rather than predicted by a model. A total of 100 Padma tomatoes of known starting stage were kept at ambient room temperature and observed over a continuous 31-day period. Each morning, between 7:00 and 7:20~a.m., every fruit was photographed from four angles, and colour differentiation was used to place it in one of the four stages from Table~\ref{tab:ripeness}; the number of days spent in each stage was recorded. A fruit left the study once it became unsellable -- noticeably soft, mouldy or rotting, or badly blemished.

\TODO{number of fruits per starting stage; calendar dates of the observation period; and the measured temperature and relative humidity of the storage environment (rather than "room temperature"), so the study can be compared against the published Padma study, which recorded 29\textdegree C and 80\,\% relative humidity}

The remaining shelf life for a detected class is the time left in that stage plus the full length of every stage after it -- for example, a fruit detected as breaker is credited with the rest of its time as breaker plus the entire turning and red stages that follow. This is an upper bound, not an exact value, since a single image cannot show how long a fruit has already been in its current stage. A fruit classified as defective is assigned zero shelf life, since it is rejected at the point of sorting. These per-class values are stored as a fixed lookup table used at the moment of classification; the full observation record is in Appendix~B.

\subsection{Complete End to End System Process}
\label{sec:m_endtoend}

The steps above combine into one process, shown in Figure~\ref{fig:workflow}. Fruit is collected and pre-sorted, photographed under controlled lighting, annotated, preprocessed and augmented, then split into three sets; the detector is trained and tested on these. In deployment, each live frame goes through the trained model, each fruit is tracked and classified as in Section~\ref{sec:m_realtime}, the result becomes a sorting command, the matching gate opens when the fruit reaches its sensor, and the classification, confidence and estimated shelf life are logged and shown on the dashboard.

\begin{figure}[H]
\centering
\TODO{Figure 3.11 -- complete end to end system workflow from fruit input to sorted output and logged record}
\caption{Complete end to end system workflow from fruit input to sorted output and logged record}
\label{fig:workflow}
\end{figure}

\subsection{Experimental Procedure}
\label{sec:m_experiments}

Five experiments were used to evaluate the system.

\textbf{Experiment 1 -- Held-out test evaluation.} Measures detection and classification performance on unseen data, using the dataset from Section~\ref{sec:m_dataset} and the procedure from Sections~\ref{sec:m_training} and~\ref{sec:m_evaluation}: precision, recall, mAP@50 and mAP@50--95, per class and overall, from a single pass over the fixed test set.

\textbf{Experiment 2 -- Architecture comparison.} Compares the deployed YOLOv8m model against a YOLO26m model trained on the same data and tested on the same held-out set, so the choice of architecture is backed by this project's own evidence. The two training runs were not set up identically -- one used a fixed batch size and the other automatic batch selection, and they stopped after different numbers of epochs -- and this difference is reported alongside the result rather than hidden.

\textbf{Experiment 3 -- Live classification validation.} Checks classification outside the offline test set, by showing fruit of known class to the live camera in the lightbox and recording what the dashboard reports. This experiment is what first exposed the lighting problem in Section~\ref{sec:annotation_qc}, and it was repeated after that was fixed. The result recorded is the proportion of fruits classified correctly.

\textbf{Experiment 4 -- Conveyor transport and end-to-end sorting.} Runs the conveyor with fruit of known class on the belt, in two parts: a mechanical run with no camera or sorting logic active, checking that transport is smooth and stable; and a full run with the whole pipeline active, checking how often each fruit is diverted at the correct gate.

\textbf{Experiment 5 -- Shelf-life observation.} The study described in Section~\ref{sec:m_shelflife}, recording how many days each fruit spends in each ripening stage.

Raw records for all five experiments are given in Appendix~B.

\subsection{System Testing Procedure}
\label{sec:m_testing}

Testing was planned in four phases: controlled testing under lightbox lighting; stress testing with varied fruit orientation, occlusion, distance and lighting; testing under realistic operating conditions; and comparing automated grading against manual grading. The resulting test cases are listed in Table~\ref{tab:testcases}, with outcomes and pass/fail status reported in Chapter~\ref{ch:results}.

\begin{table}[H]
\caption{System test cases and expected results}
\label{tab:testcases}
\centering
\begin{tabularx}{\textwidth}{l p{4.0cm} p{3.2cm} X}
\toprule
\textbf{ID} & \textbf{Test description} & \textbf{Input} & \textbf{Expected result} \\
\midrule
T-01 & Held-out test set detection and classification accuracy & 777 unseen labelled images & High precision, recall and mean average precision across all five classes \\
T-02 & Architecture comparison on the same held-out split & Two trained checkpoints & A defensible basis for the deployed architecture choice \\
T-03 & Live classification inside the lightbox & Known-class fruit, live camera & Correct class reported for each fruit presented \\
T-04 & Conveyor belt mechanical transport & Fresh fruit on the running belt & Smooth and stable transport without stalling or displacement \\
T-05 & Non-tomato object rejection & Frame containing fruit and distractor objects & Distractors not classified into any tomato class \\
T-06 & Single gate actuation bench test & Direct command to one gate, belt stationary & Corresponding servo opens and returns automatically \\
T-07 & End-to-end sorting accuracy & Known-class fruit moving on the running belt & Fruit physically diverted at the correct gate \\
T-08 & Wireless command reliability under retry & Sequenced class commands over the live link & Commands delivered exactly once despite intermittent packet loss \\
\bottomrule
\end{tabularx}
\end{table}

These cases cover detection accuracy (T-01, T-02), classification (T-03), mechanical function (T-04, T-06), and integration (T-05, T-07, T-08).

\subsection{Ethical and Safety Considerations}
\label{sec:m_ethics}

This research involved food produce and a prototype conveyor, not human or animal subjects or personal data, so no ethics board review was required. Wiring was checked before every powered test given the 24\,V and 12\,V supplies in use; the belt, rollers, motor and gate servos are pinch-point hazards, so the belt area was kept clear of hands during powered tests; fruit used for photography and conveyor testing was handled and disposed of hygienically, with no destructive testing other than the shelf-life study, where fruit is allowed to spoil by design; and images originally taken from a public dataset were used under that dataset's terms and were later removed, as described in Section~\ref{sec:annotation_qc}.

\clearpage
